\documentclass[11pt,a4paper]{article}
\usepackage[utf8]{inputenc}
\usepackage[T1]{fontenc}
\usepackage{amsmath, amssymb, amsthm}
\usepackage{geometry}
\geometry{margin=2.5cm}

\theoremstyle{definition}
\newtheorem{lemma}{Lemma}[section]
\newtheorem{definition}{Definition}[section]

\title{Lemma und Definition: Quaternion-Formfaktor Q(k) für QZG}
\author{Formale mathematische Grundlagen}
\date{\today}

\begin{document}

\maketitle

\section{Lemma: Exakte Fixierung eines Eigenwerts durch Nullraum-Bedingung}

\textbf{Setup (wie im Manuskript).} 
Sei $N=2m+1$ und
\[
A := \operatorname{diag}(t_{k-m},\dots,t_{k+m}) + \varepsilon R
\quad \text{auf}\quad \mathbb{C}^{N},
\]
und
\[
D :=\begin{pmatrix}0&A\\A^*&0\end{pmatrix}
\quad \text{auf}\quad \mathbb{C}^{N}\oplus \mathbb{C}^{N}.
\]
Sei $e_0\in\mathbb{C}^N$ der Standardbasisvektor, der die zentrale Position $k$ im Fenster markiert (also die Komponente zu $t_k$).

\begin{lemma}[Exakte Fixierung eines Eigenwerts]
Gilt
\[
R e_0=0\quad\text{und}\quad R^* e_0=0,
\]
so ist $e_0$ Eigenvektor von $A$ und $A^*$ mit Eigenwert $t_k$. Insbesondere sind
\[
\psi_\pm :=\frac{1}{\sqrt{2}}\binom{e_0}{\pm e_0}
\]
Eigenvektoren von $D$ mit
\[
D\psi_\pm=\pm t_k \psi_\pm.
\]
\end{lemma}

\textbf{Beweis (3 Zeilen).}
Aus $A=\operatorname{diag}(\cdot)+\varepsilon R$ folgt
\[
A e_0 = t_k e_0 + \varepsilon R e_0 = t_k e_0.
\]
Analog $A^* e_0 = t_k e_0$. Dann
\[
D\binom{e_0}{\pm e_0}=\binom{A(\pm e_0)}{A^* e_0}
=\binom{\pm t_k e_0}{t_k e_0}=\pm t_k \binom{e_0}{\pm e_0}.
\]
$\square$

\textbf{Kommentar im Tao-Ton:} Das ist ein \emph{linear-algebraisches Fixierungslemma} (kein ``Topologie-Schutz'' im technischen Sinn), aber als Designprinzip für $R$ sehr stark und exakt.

\section{Definition: Quaternion-Formfaktor $Q(k)$ als prüfbare, endliche Größe}

Wir benötigen eine Definition, die (i) \textbf{nur aus diskreten Daten} im Fenster stammt (Primklassen), (ii) die \textbf{Orientierung/$\pm$} abbildet, und (iii) \textbf{robust normalisiert} ist.

\subsection{Klassenabbildung}

Für jedes Index-Label $n$ (Primzahlindex) definiere die Familienklasse
\[
g(n)\in\{E,A,B,C\}
\qquad\text{durch}\qquad
p_n\bmod 12\in\{1,5,7,11\},
\]
mit $1\mapsto E$, $5\mapsto A$, $7\mapsto B$, $11\mapsto C$.

\subsection{Quaternionische Orientierung als Vorzeichenfunktion}

Fixiere einmalig eine Orientierung (z.B. $AB=+C$, $BC=+A$, $CA=+B$; Umkehrung gibt Minus). Definiere
\[
\operatorname{sgn}(X,Y;Z)\in\{-1,0,+1\}
\]
durch
\begin{itemize}
\item $=+1$, falls $XY=Z$ in der Quaternion-Tafel,
\item $=-1$, falls $YX=Z$ (also $XY=-Z$),
\item $=0$, sonst.
\end{itemize}
Damit ist ``Plus/Minus'' formal kodiert.

\subsection{Fenster und Gewichte}

Für ein Fenster $W_k:=\{k-m,\dots,k+m\}$ setze Gewichte $w_{ij}\ge 0$ (einfachste Wahl: $w_{ij}=1$ oder ein Taper $w_{ij}=e^{-|i-j|/\ell}$).

\subsection{Formfaktor}

\begin{definition}[Orientierter Quaternion-Formfaktor]
Definiere den \textbf{orientierten} Quaternion-Formfaktor
\[
\boxed{
Q(k)
:=
\frac{1}{Z_k}
\sum_{\substack{i,j\in W_k\\ i\neq j}}
w_{ij} \cdot \operatorname{sgn}\big(g(i),g(j);g(k)\big),
}
\qquad
Z_k:=\sum_{\substack{i,j\in W_k\\ i\neq j}} w_{ij}.
\]
Dann gilt automatisch $|Q(k)|\le 1$.
\end{definition}

Optional (oft sinnvoll): der \textbf{unorientierte} Resonanzgrad
\[
Q_0(k):=\frac{1}{Z_k}\sum_{i\neq j} w_{ij} \cdot \mathbf{1}\{g(i)g(j)=g(k)\},
\]
und man berichtet beide: $Q_0(k)$ misst \emph{Resonanzhäufigkeit}, $Q(k)$ misst \emph{Orientierungsüberschuss}.

\section{Minimale Testprozedur (falsifizierbar, robust)}

Das ist eine schlanke ``Methods''-Box, die ins Paper gesetzt werden kann.

\begin{enumerate}
\item \textbf{Parameter fixieren:} Wähle $m$ (Fenstergröße), Gewichte $w_{ij}$, sowie ggf. $\varepsilon$ und eine Temperatur $\beta$ für die Entropie $S(\beta)$ aus $Z(\beta)=\mathrm{Tr}(e^{-\beta D})$.

\item \textbf{Daten erzeugen:} Für $k$ in einem Intervall $[K,2K]$:
   \begin{itemize}
   \item berechne $Q(k)$ und/oder $Q_0(k)$,
   \item berechne den lokalen Gap $G(k):=p_{k+1}-p_k$,
   \item optional: berechne $S_k(\beta)$ aus dem lokalen Block-Dirac $D_k$.
   \end{itemize}

\item \textbf{Nullmodell:} Permutiere die Klassenlabels $g(i)$ innerhalb jedes Fensters (oder global), berechne $Q^{\mathrm{perm}}(k)$. Wiederhole $M$-mal.

\item \textbf{Teststatistik:} z.B. Korrelation $\mathrm{corr}(Q(k),G(k))$ oder Regression
   \[
   G(k)=a+b \cdot Q(k)+c\log p_k+\text{Fehler}.
   \]
   Vergleiche die beobachtete Statistik gegen die Permutationsverteilung $\to$ p-Wert/CI.

\item \textbf{Robustheit:} wiederhole für mehrere $(m,\beta,\varepsilon,w_{ij})$. Ein echter Effekt muss \textbf{stabil} bleiben.
\end{enumerate}

\textbf{Tao-kompatibler Claim}, wenn es klappt: \emph{``In dem getesteten Bereich zeigt $Q(k)$ eine statistisch signifikante, robuste Assoziation mit $G(k)$ relativ zu einem Permutationsnullmodell.''}

Wenn es nicht klappt, hast du trotzdem ein sauberes Ergebnis (und kannst die Metapher entschärfen).

\end{document}
